\documentclass[a4paper,10pt]{article}
\usepackage[cm]{fullpage}
\usepackage{graphicx}
\usepackage{hyperref}

\begin{document}

\begin{titlepage}
  \setlength{\parindent}{0cm}

  \includegraphics[width=200px]{Images/uob-logo-black-transparent}

  \Large
  Faculty of Engineering and Design

  \vspace{80pt}

  \LARGE
  Final Year MEng Project \\
  Project Plan and Literature Review

  \vspace{80pt}
  \textbf{Robot Football - High Level Control and Strategy Implementation}

  \vspace{10pt}
  \emph{Anthony Richards} \\
  \emph{\today}

  \vspace{80pt}
  Supervisor: \emph{Dr.~Pejman Iravani}

  \vspace{10pt}
  Assessor: \emph{TBD}
\end{titlepage}

\pagenumbering{roman}

\begin{abstract}
TODO : Abstract
\end{abstract}

\tableofcontents

\cleardoublepage

\pagenumbering{arabic}

\section{Project Background}
TODO : Background

\section{Aims and Objectives}
The aim of this project is to produce a set of algorithms that will allow a team of five robots to play a game of football against another team and win.  In order to achieve this, the following subsystems must be in place:

\begin{itemize}
\item A subsystem that will allow the robots to position themselves at an arbitrary position with a given final speed and direction of motion (within the limits of the robots).  This will allow the robots to interact with the ball to produce the motion required by the higher-level planning algorithms.

\item A subsystem that can predict the location of the ball or an opposing robot at a given time in the future.  The prediction will only need to work for the next two to three seconds to allow decisions to be made.

\item A subsystem that will monitor the behaviour of each robot and ensure that it's actions will not cause a foul (as described in the FIRA rules of play \cite{firaRules}).  For example, it should intervene if an action will cause a collision with an opponent by altering a planned route appropriately.  This will ensure that the team is compliant with the rules and does not incur any corrective action from the referee.

\item A subsystem that will combine the previous two to produce a path for a robot to allow it to interact with the ball to change the balls speed and direction of motion.  This will allow the higher level planning sub-systems to state a desired ball motion path and have a robot cause that action to happen.

\item A subsystem that will take in the current state of the system and decide on an appropriate course of action. When the ball is under the control of the current player, it will have the robots attack the opposing goal and attempt to score a goal.  When the ball is under the control of the opposition it will attempt to defend its own goal and disrupt the opponents actions. When the ball is not under any defined control, it will attempt to take control of the ball.
\end{itemize}

The entire system must be sufficiently lightweight to allow all of it to be run once per simulator frame.  This will allow the system to respond quickly to the current state of the system and take advantage of any changes as soon as possible.

A successful system will be able to consistently perform better than the standard strategy provided with the SimuroSot simulator.  This can be measured by setting the two strategies against each other in the simulator and seeing which team wins.

If a successful system is produced, it may be possible to present it to the FIRA competition happening this year in Bristol.  This will be dependent on the progress made, as the project will have restricted man-hours available for work compared to other teams in the SimuroSot league (as other teams usually consist of more than just one person).

\section{Methodology and Deliverables}
The main result of the project will be a SimuroSot simulator compatible DLL containing the algorithms required to produce the desired strategy.  This DLL can then be loaded by the simulator and the simulation run to observe the results.

As the documentation for the simulator is sparse, it will first be necessary to examine the behaviour of the system to confirm the underlying equations of motion and what (if any) simplifications have been made to the physics.  This information can then be used to produce the motion prediction algorithms, as well as the algorithms required to plan trajectories.

The initial attempt at high-level strategy will use a scoring system to determine which robots are in the best position to achieve the current goal (e.g. offence or ball interception).  These scores can then be used to produce a tree of possible outcomes, and a suitable algorithm can be used to select the best move.  This will attempt to take into account not just the current 'move', but also the next one, to best allow the system to plan ahead.

The system will be made as lightweight as possible by taking full advantage of the advanced features of modern processors that allow multiple actions to be achieved simultaneously (e.g. with multi-threading and Single-Instruction Multiple-Date instructions as discussed in \cite{intelMatrixInverse}).  The profiling tools provided with Visual Studio will be used to confirm any perceived improvements.

The success of the system will be measured by comparing the behaviour of the system over a number of games compared to the basic strategy playing against itself.  The best strategy would always beat the standard strategy, but one that creates a significant improvement on the chances of success would still be a success.

\section{Resources}
As this project is targeted at the SimuroSot league, all the activities will be carried out in software only.  In order to produce the required software, the following software will be required:

\begin{itemize}
 \item Microsoft Visual Studio 2010
 \item SimuroSot Middle League Simulator
\end{itemize}

Visual Studio is available on the University computers in the Maxwell and Faraday labs in the Electrical Engineering Department, as well as on the GigaTerms server which is available via Remote Desktop.  It is also available free of charge to students from the Microsoft Dreamspark Website \cite{dreamsparkVS}. The SimuroSot Simulator is also available free of charge from the FIRA website \cite{simurosotSim}.

In order to run the required software, access to a PC running Microsoft Windows will be required.  These computers are readily available in the Maxwell laboratory in the Electrical Engineering Department, as well as the Software Projects Lab in the same department.  A personal computer capable of running the required software is also available, should it be required.

The development of the project will be aided by the use of an online note-taking website (EverNote), and source code control will be achieved using Subversion.  Combined, these will allow development of the project to be carried on any computer with the suitable software, removing the need for dedicated lab space for the project.  All the required support software is readily available on the University computers, as well as the personal machine.

\section{Predicted Timescale}
TODO : Gannt Chart

\cleardoublepage
\bibliographystyle{plain}
\phantomsection
\addcontentsline{toc}{section}{References}
\bibliography{References/references}

\end{document}
